\chapter{Checkpoint Semantics Reference}
\label{app:checkpoints}

This appendix collects, in full, the checkpoint semantics introduced in
\Cref{ch:lessons-05-07}, for use as a standing reference when designing
resumable, long-running scientific workflows.

\section{Checkpointing as a Special Case of the Provenance Graph}

Before working through the mechanics, it is worth connecting checkpoint
semantics explicitly to the provenance-graph concept developed in
\Cref{ch:schema}. A checkpoint file, in the precise sense developed there,
is a \emph{partial, persisted snapshot} of the provenance graph's node
annotations --- specifically, the \texttt{working\_dir} and final exit
\texttt{code} of every task that has completed at least once --- written to
a location (the checkpoint JSON file) that survives independently of the
Python process that produced it. This framing clarifies precisely what a
checkpoint file does and does not capture: it captures exactly the same two
facts \Cref{ch:lessons-03-04}'s debugging procedure already treats as
foundational --- where a task's output lives, and whether the task
succeeded --- and nothing more. It does not, in particular, capture the
\texttt{launcher.sh} content itself, which is why, as the limitations
section below reiterates, checkpoint reuse cannot by itself detect that a
task's \emph{command} changed between the run that produced the checkpoint
entry and the run attempting to reuse it.

\section{What Checkpointing Records}

When a \texttt{Workflow} is created with \texttt{checkpoint\_file}, DAGonStar
writes a JSON file containing: the workflow scratch base; per-task working
directories; workflow name; task name; and the task exit code after
execution.

\begin{lstlisting}[caption={Representative checkpoint JSON shape.}]
{
  "_scratch_dir": "/tmp/",
  "Experiment.A": {
    "working_dir": "/tmp/1710000000000-A",
    "workflow": "Experiment",
    "name": "A",
    "code": 0
  }
}
\end{lstlisting}

\section{Creating and Resuming a Checkpointed Workflow}

\begin{lstlisting}[caption={Creation.}]
from dagon import Workflow
from dagon.task import DagonTask, TaskType


workflow = Workflow("Experiment", checkpoint_file="checkpoint.json")

task = DagonTask(TaskType.BATCH, "A", "echo result > result.txt")
workflow.add_task(task)
workflow.make_dependencies()
workflow.run()
\end{lstlisting}

\begin{lstlisting}[caption={Resumption.}]
workflow = Workflow("Experiment", checkpoint_file="checkpoint.json")
task = DagonTask(TaskType.BATCH, "A", "echo result > result.txt")
workflow.add_task(task)
workflow.make_dependencies()
workflow.run(resume_checkpoint_file="checkpoint.json")
\end{lstlisting}

If the checkpoint records \texttt{code: 0} for \texttt{Experiment.A} and the
recorded working directory still exists, the task can be treated as already
complete and is skipped on the resumed run.

\section{Task Identity}

Checkpoint keys are built from \texttt{<workflow-name>.<task-name>}. Stable
workflow and task names are therefore essential: renaming a task is, from
the checkpoint system's perspective, indistinguishable from introducing an
entirely new task with no recorded history.

This naming-stability requirement deserves a slightly fuller treatment than
a single warning sentence, since it interacts with ordinary software
maintenance in a way that is easy to overlook. As a workflow's Python source
evolves over the lifetime of a research project --- task names refactored
for clarity, a workflow renamed to reflect a renamed project --- every such
rename silently orphans the corresponding checkpoint entries: the old key
remains in the checkpoint JSON, referring to a working directory that may
still exist on disk, while the renamed task generates a new key with no
checkpoint history of its own. This is not a bug; it follows directly from
the dictionary-key-based identity scheme described above. It does mean,
however, that a workflow author renaming tasks in an actively used,
checkpointed workflow should treat the rename as equivalent, from a
checkpoint-continuity perspective, to discarding all prior checkpoint
history for the renamed tasks, and plan accordingly --- for instance, by
completing any in-flight, expensive run under the old names before applying
the rename, rather than renaming mid-run and hoping for seamless resumption.

\section{Working Directory Persistence}

Checkpoint reuse requires the recorded working directory to exist. If
scratch directories are deleted, moved, or garbage-collected, checkpoint
reuse cannot recover those task outputs. For long-running scientific
workflows, scratch directories should be placed on storage with a retention
policy appropriate to the experiment's expected lifetime.

\section{Checkpointing and Dataflow}

Checkpointing is most valuable for dataflow workflows whose upstream tasks
produce large intermediate files. Given the chain

\begin{lstlisting}[basicstyle=\ttfamily\small, frame=none]
download -> preprocess -> simulate -> postprocess
\end{lstlisting}

if \texttt{simulate} fails, a later run may reuse \texttt{download} and
\texttt{preprocess} if their working directories and checkpoint entries
remain valid --- avoiding, potentially, hours of redundant data acquisition
and preprocessing for a single failed simulation attempt.

This four-stage example is worth re-examining quantitatively, building on
\Cref{ch:lessons-05-07}'s staging-cost discussion: the value checkpointing
delivers in this scenario is directly proportional to the cost of the
stages it allows a researcher to \emph{skip} on resumption, relative to the
cost of the stage that actually failed. If \texttt{download} and
\texttt{preprocess} together consume four hours of wall-clock time and
\texttt{simulate} consumes thirty minutes before failing, checkpoint reuse
converts a worst-case four-and-a-half-hour recovery cost into a thirty-
minute one --- a substantial, quantifiable return on the comparatively
small overhead of supplying a \texttt{checkpoint\_file} argument. The same
calculation, run in reverse, also clarifies when checkpointing delivers
little practical benefit: a workflow whose every stage is roughly equally
cheap gains comparatively little from checkpoint reuse, since even a full
restart from scratch is itself inexpensive.

\section{Explicit Checkpoint Tasks}

\texttt{TaskType.CHECKPOINT} maps to \texttt{dagon.checkpoint.Checkpoint}.
This task type is part of the implementation for explicit checkpoint
behavior; the workflow-level \texttt{checkpoint\_file} mechanism described
above remains the primary, user-facing mechanism for recording and resuming
task metadata.

\section{Remote, Docker, and Cloud Checkpointing}

Remote checkpoint reuse checks for the existence of the recorded directory
on the remote host through the task's SSH connection, which means:

\begin{itemize}[leftmargin=1.4em]
  \item the same host must be reachable;
  \item credentials must still work;
  \item the remote scratch directory must still exist;
  \item paths must be valid on the remote filesystem.
\end{itemize}

\begin{cautionbox}
Docker and cloud task checkpointing depends on the persistence of the host
scratch directory and, for cloud tasks, on the persistence or
reconstruction of the remote execution environment. Do not assume a
terminated cloud instance preserves local scratch data unless it was
explicitly written to persistent storage.
\end{cautionbox}

These three additional preconditions --- host reachability, credential
validity, and remote filesystem persistence --- compose directly with the
backend-specific operational considerations already developed in
\Cref{ch:backends}: a remote or Slurm-scheduled task's queueing and latency
costs (already discussed there) determine how \emph{expensive} a failure
partway through that task is, while this section's preconditions determine
whether checkpoint reuse can actually \emph{recover} from that failure
without re-incurring those costs. A workflow author relying on checkpoint
reuse for a remote or cloud-backed stage should verify both sets of
considerations together, not in isolation.

\section{Scientific Usage Pattern}

For expensive workflows, the documentation recommends the following
seven-step discipline:

\begin{enumerate}[leftmargin=1.6em]
  \item Use a stable workflow name.
  \item Use stable task names.
  \item Set \texttt{checkpoint\_file}.
  \item Set a scratch base with enough storage.
  \item Preserve the checkpoint JSON with experiment metadata.
  \item Preserve or archive task working directories needed for reuse.
  \item On rerun, use \texttt{resume\_checkpoint\_file}.
\end{enumerate}

\section{Limitations and Cautions}

\begin{cautionbox}
\begin{itemize}[leftmargin=1.4em]
  \item Checkpoints record task \emph{success} and working directory, not
    full semantic equivalence of code, inputs, or parameters.
  \item If a command changes but the task name and previous success remain,
    checkpoint reuse may incorrectly skip recomputation.
  \item For rigorous scientific reproducibility, store code version,
    configuration, input checksums, and software environment \emph{alongside}
    the checkpoint file, not only inside it.
  \item Garbage collection or manual cleanup can invalidate checkpoint
    reuse.
\end{itemize}
\end{cautionbox}

The first two items of this list deserve to be read together, since they
describe the same underlying limitation from two angles: a checkpoint
entry's identity is purely nominal (task name) and purely behavioral
(recorded exit code), never \emph{content-based}. A content-based scheme
--- recording, say, a hash of the task's fully resolved command, as
captured in \texttt{launcher.sh} --- would in principle detect a changed
command automatically and correctly invalidate the corresponding checkpoint
entry, at the cost of additional implementation complexity the current,
simpler name-and-status scheme avoids. A workflow author who needs this
stronger guarantee today must provide it manually, by following the fifth
bullet's recommendation to preserve code version and configuration
alongside the checkpoint file, and by manually deleting or ignoring stale
checkpoint entries whenever a task's command is deliberately changed.

\section{Verification}

\begin{lstlisting}[style=bash]
cd examples/dataflow/batch
python dataflow-checkpoint-demo.py
\end{lstlisting}

Then inspect the generated checkpoint file named by the example script. A
minimal custom test illustrating two successive resumed runs:

\begin{lstlisting}
from dagon import Workflow
from dagon.task import DagonTask, TaskType


def build():
    workflow = Workflow("CheckpointDemo", checkpoint_file="checkpoint.json")
    task = DagonTask(TaskType.BATCH, "A", "echo data > data.txt")
    workflow.add_task(task)
    workflow.make_dependencies()
    return workflow


build().run(resume_checkpoint_file="checkpoint.json")
build().run(resume_checkpoint_file="checkpoint.json")
\end{lstlisting}

The second run can reuse the completed task if the recorded directory
remains on disk.
